\documentclass[conference]{IEEEtran}

\usepackage{cite}
\usepackage{graphicx}
\usepackage{amsmath,amssymb}
\usepackage{booktabs}
\usepackage{multirow}
\usepackage{url}
\usepackage{xcolor}
\usepackage{balance}

\graphicspath{{figures/}}

\begin{document}

\title{MSJC: Multi-Stage Jamming Classification for\\5G O-RAN with Closed-Loop MLOps}

\author{
\IEEEauthorblockN{Hyungmin Kim and Eric W. Burger}
\IEEEauthorblockA{Virginia Tech, Arlington, VA, USA\\
\{henrykim9319, ewburger\}@vt.edu}
}

\maketitle

% ================================================================
\begin{abstract}
Jamming attacks threaten 5G networks, yet most existing detection methods rely on simulated data and lack deployment on standardized O-RAN architectures. We present MSJC, a Multi-Stage Jamming Classification framework deployed on a real 5G O-RAN testbed that detects seven distinct jamming attack types---including protocol-aware attacks targeting PSS, PDCCH, and DMRS---with 100\% detection for six of seven modes (91.9\% for PDCCH) and 5.3\% false alarm rate. Our three-stage cascade combines: (1)~an MLP for fast 5-class KPI-based classification, (2)~a sliding-window KSVM that extracts 12-dimensional temporal statistics to detect intermittent attacks invisible to per-sample classifiers, and (3)~a MobileNetV3 for I/Q spectrogram-based protocol-aware sub-classification achieving 85\% accuracy across four sub-types. We further implement a closed-loop MLOps pipeline with GMM-based auto-labeling and automated retraining via the A1 interface. All results are validated via blocked temporal 5-fold cross-validation on 1,492 real testbed samples with zero sample overlap between train and test sets. Experiments on the CCI~xG Testbed with three USRP~X310 software-defined radios demonstrate end-to-end inference latency $\leq$23\,ms (p95), well within the Near-RT~RIC 100\,ms budget.
\end{abstract}

\begin{IEEEkeywords}
5G, O-RAN, jamming detection, xApp, machine learning, Near-RT RIC, E2SM-KPM, closed-loop MLOps
\end{IEEEkeywords}

% ================================================================
\section{Introduction}
\label{sec:intro}

The Open Radio Access Network (O-RAN) architecture disaggregates the traditional RAN into programmable, multi-vendor components controlled by RAN Intelligent Controllers (RICs)~\cite{polese2023, abdalla2022}. While this openness enables innovation through xApps and rApps~\cite{santos2025, hoffmann2024}, it also expands the attack surface, making 5G networks vulnerable to RF jamming~\cite{pirayesh2022, harvanek2024}. Tactical and military 5G deployments in contested RF environments face this threat most directly, as adversaries may employ protocol-aware electronic warfare (EW) techniques.

Modern jamming attacks against 5G New Radio (NR) are no longer limited to broadband noise. Protocol-aware jammers can selectively target synchronization signals (PSS/SSS), control channels (PDCCH), or reference signals (DMRS), each producing distinct spectral signatures but minimal impact on aggregate KPI metrics~\cite{savadatti2024, shahid2024}. Detecting and classifying these attack types in real time is necessary for timely mitigation.

Existing ML-based jamming detection approaches have two main limitations. First, most rely on simulated data without validation on real 5G hardware~\cite{hachimi2020, kryszkiewicz2023}. Jere et al.~\cite{jere2024} demonstrated Bayesian-assisted ML on a 5G NR testbed but did not integrate with O-RAN interfaces. Second, static models cannot adapt to evolving attack patterns; Rahman et al.~\cite{rahman2025} proposed closed-loop retraining but limited to binary classification in simulation. Recent O-RAN-based approaches have begun addressing jamming on real testbeds~\cite{moore2025, bogucka2025, pazperez2025}, but none combine multi-class classification with protocol-aware sub-typing and automated retraining in a single framework.

In this paper, we present MSJC (Multi-Stage Jamming Classification), a framework implemented as a Near-RT RIC xApp for real-time inference with Non-RT RIC rApps for closed-loop MLOps, validated on the CCI~xG Testbed~\cite{dasilva2025}. Our contributions are:

\begin{enumerate}
    \item \textbf{3-Stage Classification Cascade:} Extending Hachimi et al.'s two-stage MLP$+$KSVM design~\cite{hachimi2020} from C-RAN to O-RAN, we add a third stage (MobileNetV3~\cite{howard2019}) for protocol-aware spectrogram classification, redesign the KSVM with sliding-window temporal features, expand from 4 to 7 jamming modes, and validate on real 5G hardware---achieving 100\% detection for six modes with 5.3\% false alarm rate and $\leq$23\,ms latency (p95), evaluated via blocked temporal cross-validation.

    \item \textbf{Sliding-Window Statistical Features:} A 12-dimensional feature vector extracted from 15-second KPI windows that captures temporal patterns (BLER variance, periodicity, state transitions) invisible to per-sample classifiers, enabling detection of intermittent Random and Reactive jamming.

    \item \textbf{Closed-Loop MLOps on Real O-RAN:} Adopting Rahman et al.'s SAJD architecture~\cite{rahman2025}, we implement GMM-based Labeler rApp and Training Manager rApp and extend them from binary detection to multi-class classification, with end-to-end validation on a real 5G testbed---the first such demonstration on actual hardware.
\end{enumerate}


% ================================================================
\section{Related Work}
\label{sec:related}

Pirayesh and Zeng~\cite{pirayesh2022} survey jamming attacks and countermeasures across wireless networks; Savadatti et al.~\cite{savadatti2024} and Shahid et al.~\cite{shahid2024} provide 5G-specific taxonomies. Harvanek et al.~\cite{harvanek2024} survey PHY-layer threats including jamming, noting the growing role of ML-based detection.

Hachimi et al.~\cite{hachimi2020} proposed a two-stage MLP$+$KSVM cascade for 4-class jamming classification in 5G C-RAN, achieving 94.51\% on simulated data but without real RF validation, O-RAN integration, or temporal feature extraction. Jere et al.~\cite{jere2024} demonstrated Bayesian ML on a real 5G NR testbed (AUC~$>$0.988, 180\,ms delay) but without O-RAN interfaces or protocol-aware classification. Our work extends Hachimi's cascade to three stages (adding MobileNetV3~\cite{howard2019} for spectrogram-based protocol sub-typing), redesigns the KSVM with sliding-window temporal statistics, expands to seven jamming modes, and deploys on a real O-RAN testbed.

Several works deploy security xApps on the Near-RT RIC: Moore et al.~\cite{moore2025} demonstrated LSTM-based intrusion detection on a real testbed; Bogucka et al.~\cite{bogucka2025} implemented jamming detection with MCS adaptation; Paz-P\'{e}rez et al.~\cite{pazperez2025} designed a four-xApp cascade on FlexRIC; and Reus-Muns et al.~\cite{reusmuns2024} validated spectrogram-based detection via a YOLO xApp. However, none combine multi-class classification with protocol-aware sub-typing and closed-loop retraining. Rahman et al.~\cite{rahman2025} proposed SAJD with GMM-based auto-labeling and ClearML-driven retraining via the A1 interface, but limited to binary detection in simulation. We adopt Rahman's MLOps architecture and extend it to 7-class classification on real hardware.


% ================================================================
\section{System Architecture}
\label{sec:architecture}

\subsection{CCI xG Testbed}

All experiments are conducted on the CCI~xG Testbed~\cite{dasilva2025}, an O-RAN-compliant research platform at Virginia Tech. Our setup uses three instances:

\begin{itemize}
    \item \textbf{Instance-1 (UE):} srsUE running 5G NR SA on Band~3 FDD (1842.5\,MHz DL, 20\,MHz BW, 15\,kHz SCS, 106~PRBs).
    \item \textbf{Instance-2 (gNB + RIC):} srsRAN Project gNB with E2SM-KPM support, Open5GS 5G Core, and FlexRIC Near-RT RIC hosting the MSJC xApp.
    \item \textbf{Instance-3 (Jammer):} A custom 7-mode jammer generating protocol-aware RF interference, with OFDM structure matched to the gNB (FFT=1536, CP=108).
\end{itemize}

Each instance is equipped with an Ettus USRP X310 SDR connected via 10G Ethernet (MTU~9000). The gNB reports per-UE KPI metrics via E2SM-KPM at 1-second intervals. Fig.~\ref{fig:arch} shows the overall system architecture.

\begin{figure}[t]
\centering
\includegraphics[width=\columnwidth]{fig_architecture.pdf}
\caption{MSJC testbed architecture on the CCI~xG Testbed. The MSJC xApp runs on the FlexRIC Near-RT RIC, receiving E2SM-KPM telemetry from the gNB. Non-RT RIC rApps automate labeling and retraining via the A1 interface.}
\label{fig:arch}
\end{figure}

\subsection{Jammer Design}

Our custom jammer generates seven distinct attack modes, with OFDM symbol structure (FFT=1536, CP=108, 14~symbols/slot) matched to the target gNB to maximize interference effectiveness:

\begin{itemize}
    \item \textbf{Constant:} Continuous broadband AWGN (amplitude~0.5, gain~0\,dB). Produces sustained BLER elevation ($>$0.7) and SINR degradation ($-$2.2\,dB).
    \item \textbf{Random:} Frame-based random-duty-cycle pulses (30\% duty, amplitude~0.6). Creates intermittent BLER spikes with high variance.
    \item \textbf{Reactive:} Slot-aware adaptive pulses (40\% duty, amplitude~0.6). Produces periodic BLER patterns detectable through autocorrelation.
    \item \textbf{Deceptive:} OFDM-like waveform mimicking legitimate transmissions (amplitude~0.6). Causes minimal KPI degradation (BLER $\approx$ 0).
    \item \textbf{PSS/SSS:} Narrowband attack on synchronization signals (6--12 center RBs, gain~10\,dB). Targets initial cell access.
    \item \textbf{PDCCH:} CORESET-targeted attack on control channels (slots~0--3 OFDM symbols, gain~5\,dB). Disrupts scheduling information.
    \item \textbf{DMRS:} Comb-pattern attack on demodulation reference signals (subcarrier spacing~4--6, gain~10\,dB). Corrupts channel estimation.
\end{itemize}

\subsection{MSJC Classification Pipeline}

The xApp implements a three-stage cascade optimized for the Near-RT RIC's 100\,ms latency budget:

\textbf{Stage~1 (MLP):} Classifies each 8-dimensional KPI vector into five classes (Normal, Constant, Random, Reactive, Deceptive) in 1--2\,ms. Constant, Random, and Reactive detections are confirmed immediately.

\textbf{Stage~2 (KSVM):} Re-examines Normal verdicts using a sliding window of $W=15$ samples (15~seconds). Twelve statistical features capture temporal patterns invisible to per-sample analysis. An RBF-kernel SVM decides Normal vs.\ Attack in 3--4\,ms.

\textbf{Stage~3 (MobileNetV3):} Invoked only for Deceptive detections or Stage~2 false-negative catches. Raw I/Q snapshots (16,384 complex samples via passive USRP reception) are converted to 224$\times$224 spectrograms~\cite{reusmuns2024} and classified into PSS/SSS, PDCCH, DMRS, or Generic Deceptive in 8--21\,ms.

\noindent Fig.~\ref{fig:pipeline} illustrates the complete classification cascade and decision flow.

\begin{figure}[t]
\centering
\includegraphics[width=\columnwidth]{fig_pipeline.pdf}
\caption{MSJC three-stage classification cascade. Stage~1 classifies per-sample KPIs; Stage~2 re-examines Normal verdicts via sliding-window statistics; Stage~3 performs protocol-aware I/Q spectrogram classification. All paths complete within 23\,ms (p95).}
\label{fig:pipeline}
\end{figure}

\subsection{Closed-Loop MLOps}

Two Non-RT RIC rApps automate the ML lifecycle, adopting the SAJD closed-loop architecture proposed by Rahman et al.~\cite{rahman2025} and extending it from binary detection to multi-class classification:
\begin{itemize}
    \item \textbf{Labeler rApp:} Implements GMM-based unsupervised labeling pipeline (moving average smoothing $\rightarrow$ ARC drift detection with $\tau = 0.0004$ $\rightarrow$ 2-component GMM clustering) to generate binary Normal/Attack annotations from raw KPI telemetry.
    \item \textbf{Training Manager rApp:} Monitors xApp accuracy via InfluxDB, triggers Stage~1 MLP retraining when accuracy drops below 30\%, and deploys updated models via an A1-like HTTP interface---extending Rahman's ClearML-based orchestration with a standalone fallback mode.
\end{itemize}


% ================================================================
\section{Methodology}
\label{sec:methodology}

\subsection{KPI Feature Extraction}

Each E2SM-KPM indication message from the gNB is parsed and mapped to an 8-dimensional feature vector $\mathbf{x} \in \mathbb{R}^8$:

\begin{equation}
\mathbf{x} = [\text{RSRP}, \text{RSRQ}, \text{SINR}, \text{BLER}, \text{NACK}, \text{Tput}, \text{CQI}, \text{HARQ}]
\end{equation}

\noindent where RSRP uses PUSCH SNR as a proxy, SINR is derived from PUCCH measurements, and CQI is automatically scaled from the E2SM-KPM range [0--15] to the accumulator range [0--50] used during training (a $\times$3.33 correction factor identified during real testbed integration). The current srsRAN E2SM-KPM implementation exposes five directly measured KPIs (PUSCH~SNR, PUCCH~SINR, DL~BLER, NACK~ratio, CQI); RSRQ, throughput, and HARQ retransmission ratio are reserved for commercial gNB deployments that report the full O-RAN KPM measurement set. All reported results therefore reflect classification performance using five active feature dimensions---a conservative baseline that full KPM support would only improve. In practice, SINR, BLER, and CQI serve as the primary discriminators: Constant jamming causes SINR to drop from $+$5.6 to $-$2.2\,dB and BLER to spike to 0.75, while protocol-aware attacks (PSS, PDCCH, DMRS) leave all KPIs virtually unchanged (Table~\ref{tab:kpm_stats}).

\subsection{Stage 1: MLP 5-Class Classifier}

Following Hachimi et al.'s two-stage design~\cite{hachimi2020}, our Stage~1 uses an MLP as the first-pass classifier. We extend the original 4-class scheme to 5~classes (adding Deceptive as a catch-all for protocol-aware attacks) and replace raw I/Q features with 8-dimensional E2SM-KPM metrics. The architecture (8$\rightarrow$256$\rightarrow$128$\rightarrow$64$\rightarrow$5) with BatchNorm, ReLU, and Dropout~(0.3) totals 44,677 parameters. Training uses a hybrid dataset: synthetic KPI profiles calibrated to real testbed measurements (500 per class) mixed with real labeled samples from the testbed (approximately 765 real samples per fold after filtering). Clean-looking attack samples (BLER~$<$~0.03 and SINR~$>$~3\,dB and CQI~$>$~48) are deliberately excluded from Stage~1 training, as these protocol-aware attacks are designed to be caught by Stage~2 and classified by Stage~3. To prevent train/test leakage, both stages are evaluated via blocked temporal 5-fold cross-validation (detailed in Section~\ref{sec:results}).

\subsection{Stage 2: Sliding-Window KSVM}

Stage~1 classifies each KPI sample independently, which is effective for attacks causing persistent degradation (e.g., Constant) but fails to detect protocol-aware attacks that leave per-sample KPIs within normal ranges. To address this, Stage~2 re-examines all Stage~1 ``Normal'' verdicts by analyzing the \textit{temporal statistics} of a sliding window of $W=15$ consecutive KPI reports (15~seconds at 1\,Hz reporting rate).

We extract 12 statistical features from each window:

\begin{equation}
\mathbf{f} = [\mu_b, \sigma^2_b, b_{\max}, r_b, \mu_s, \sigma^2_s, s_{\min}, \sigma^2_c, c_{\min}, \mu_n, \rho_b, \tau_b]
\end{equation}

\noindent where $b$, $s$, $c$, $n$ denote BLER, SINR, CQI, NACK respectively. The features are designed to capture distinct temporal signatures:

\begin{itemize}
    \item \textbf{BLER statistics} ($\mu_b$, $\sigma^2_b$, $b_{\max}$, $r_b$): Mean, variance, maximum, and spike ratio (fraction of samples with BLER~$>$~0.03). Constant jamming produces high $\mu_b$ ($>$0.35); Random/Reactive produce intermittent BLER spikes with high variance and transitions.
    \item \textbf{SINR/CQI statistics} ($\mu_s$, $\sigma^2_s$, $s_{\min}$, $\sigma^2_c$, $c_{\min}$): Capture signal quality degradation and variability. Protocol-aware attacks produce a distinct SINR signature ($\mu_s \approx 4.7$\,dB) compared to Normal traffic ($\mu_s \approx 5.6$\,dB), despite leaving BLER near zero. CQI variance ($\sigma^2_c$) further separates Normal traffic ($\sigma^2_c \approx 11$) from protocol-aware attacks ($\sigma^2_c \approx 1.2$), reflecting the natural scheduling-driven CQI fluctuations absent under attack.
    \item \textbf{Temporal dynamics} ($\rho_b$, $\tau_b$): BLER autocorrelation at lag-1 distinguishes Reactive (periodic, high $\rho_b$) from Random (aperiodic, low $\rho_b$) jamming; state transition count $\tau_b$ captures BLER on/off switching frequency.
\end{itemize}

While Hachimi et al.~\cite{hachimi2020} used a per-sample KSVM as a binary false-negative recheck, our sliding-window approach captures temporal dynamics invisible to single-sample analysis. The RBF-kernel SVM ($C$=0.5, $\gamma$=scale) is trained on a hybrid dataset combining real testbed sliding windows with synthetic windows (500 per class). A lower regularization ($C$=0.5 vs.\ the typical $C$=10) reduces overfitting to the specific SINR distribution of the training blocks, improving generalization across temporal non-stationarity.

\subsection{Stage 3: MobileNetV3 Protocol Classification}

Stage~3 is invoked only when Stage~1 detects a Deceptive attack or Stage~2 catches a false negative, keeping its overhead off the critical path for most samples. Raw I/Q snapshots (16,384 complex samples) are captured via passive USRP reception on the UE instance and converted to 224$\times$224 spectrograms using Welch's method (segment length 256, 78\% overlap), with power spectral density mapped to normalized dB scale.

MobileNetV3-Small~\cite{howard2019} (1.52M parameters) classifies these spectrograms into four protocol-aware sub-types: PSS/SSS (center-concentrated energy in 6--12~RBs), PDCCH (time-domain CORESET bursts in first 1--3~OFDM symbols), DMRS (comb-pattern at subcarrier intervals of 4--6), and Generic Deceptive (uniform broadband OFDM). Initial training uses 1,200 synthetic spectrograms (300 per class) generated from frequency-domain attack models, followed by fine-tuning with 200 real I/Q captures (50 per class) collected via full-duplex USRP~X310 passive reception during active jamming, reducing the sim-to-real gap as motivated by SenseORAN~\cite{reusmuns2024}.


% ================================================================
\section{Experimental Results}
\label{sec:results}

\subsection{Dataset and Configuration}

We collected 1,492 KPI samples (1-second intervals) across all seven jamming modes plus Normal baseline on the CCI~xG Testbed~\cite{dasilva2025}: 678~Normal, 100~Constant, and 119~each for Random, Reactive, Deceptive, PSS, PDCCH, and DMRS. Additionally, 400 I/Q snapshots (50 per mode $\times$ 8 modes) were captured via passive USRP~X310 reception during active jamming. Table~\ref{tab:kpm_stats} summarizes the measured KPI statistics.

\begin{table}[t]
\centering
\caption{Measured KPI Statistics per Jamming Mode}
\label{tab:kpm_stats}
\small
\setlength{\tabcolsep}{3pt}
\begin{tabular}{lrrrr}
\toprule
Mode & N & SINR (dB) & BLER & CQI \\
\midrule
Normal & 678 & $+5.6 \pm 2.5$ & $0.032 \pm 0.16$ & $48.6 \pm 6.1$ \\
Constant & 100 & $-2.2 \pm 6.5$ & $0.749 \pm 0.34$ & $34.8 \pm 21$ \\
Random & 119 & $+4.2 \pm 0.9$ & $0.328 \pm 0.25$ & $49.2 \pm 2.0$ \\
Reactive & 119 & $+4.2 \pm 0.9$ & $0.275 \pm 0.25$ & $49.4 \pm 1.7$ \\
Deceptive & 119 & $+4.7 \pm 0.6$ & $0.000 \pm 0.00$ & $49.8 \pm 1.2$ \\
PSS & 119 & $+4.8 \pm 0.6$ & $0.008 \pm 0.06$ & $49.8 \pm 1.2$ \\
PDCCH & 119 & $+4.7 \pm 0.6$ & $0.007 \pm 0.08$ & $49.9 \pm 0.9$ \\
DMRS & 119 & $+4.7 \pm 0.6$ & $0.006 \pm 0.07$ & $49.8 \pm 1.0$ \\
\bottomrule
\end{tabular}
\end{table}

As illustrated in Fig.~\ref{fig:kpm_ts}, protocol-aware attacks leave per-sample KPIs virtually unchanged, but produce distinct multi-dimensional temporal signatures detectable through sliding-window analysis (Section~\ref{sec:methodology}). Both stages use \textit{blocked temporal 5-fold cross-validation}: each mode's time-series is split into five consecutive blocks, and each fold trains on four blocks (augmented with synthetic data) and tests on the held-out block, ensuring zero sample overlap between train and test sets.

\subsection{Detection Performance}

Table~\ref{tab:detection} compares Stage~1 alone versus the combined pipeline. Stage~1 detects Constant jamming (97.0\%) through sustained BLER elevation, and partially detects Random (68.9\%) and Reactive (65.6\%) through intermittent BLER spikes. However, protocol-aware attacks evade per-sample classification entirely---Stage~1 detects $<$1\% of these four modes (Table~\ref{tab:kpm_stats}).

The sliding-window KSVM in Stage~2 recovers nearly all false negatives by exploiting joint differences in SINR and CQI temporal statistics, as described in Section~\ref{sec:methodology}. The combined pipeline achieves \textbf{100\% detection for six of seven modes}, with PDCCH at 91.9\% (limited by one temporal block where the narrow CORESET-targeted interference falls below the SVM's decision boundary), and 5.3\% combined false alarm rate.

\begin{table}[t]
\centering
\caption{Detection Rate: Stage~1 vs.\ Stage~2 vs.\ Combined, evaluated via blocked temporal 5-fold CV ($N$=1,492). $\pm$ values denote standard deviation across folds.}
\label{tab:detection}
\small
\begin{tabular}{lrrr}
\toprule
Mode & Stage~1 & Stage~2 & Combined \\
\midrule
Normal (FA) & 3.7$\pm$4.8\% & 1.6$\pm$3.2\% & $\sim$\textbf{5.3\%} \\
Constant & 97.0$\pm$6.0\% & 100\% & \textbf{100\%} \\
Random & 68.9$\pm$7.4\% & 100\% & \textbf{100\%} \\
Reactive & 65.6$\pm$6.2\% & 100\% & \textbf{100\%} \\
Deceptive & $<$1\% & 100\% & \textbf{100\%} \\
PSS & 0\% & 100\% & \textbf{100\%} \\
PDCCH & $<$1\% & 91.8$\pm$12.3\% & \textbf{91.9\%} \\
DMRS & 0\% & 100\% & \textbf{100\%} \\
\bottomrule
\end{tabular}
\end{table}


\begin{figure}[t]
\centering
\includegraphics[width=\columnwidth]{fig_kpm_timeseries.pdf}
\caption{KPI time-series under different jamming modes. Constant jamming causes sustained BLER elevation; Random/Reactive produce intermittent spikes; protocol-aware attacks (PSS, PDCCH, DMRS) remain invisible to per-sample KPI analysis.}
\label{fig:kpm_ts}
\end{figure}

\subsection{Protocol-Aware Classification}

Stage~3 MobileNetV3, fine-tuned with real I/Q captures from the testbed, achieves 94\% for PSS/SSS, 72\% for PDCCH, 84\% for DMRS, and 90\% for Generic Deceptive (85\% overall, evaluated on 200 held-out real I/Q samples). Fig.~\ref{fig:stage3} shows per-class accuracy on real testbed captures. Fig.~\ref{fig:spectrograms} shows representative spectrograms from real I/Q captures, illustrating the distinct spectral signatures that enable protocol-aware sub-typing: PSS/SSS concentrates energy in central subcarriers, PDCCH shows CORESET-aligned time-domain bursts, and DMRS exhibits a characteristic comb pattern.

\begin{figure}[t]
\centering
\includegraphics[width=\columnwidth]{fig_stage3_accuracy.pdf}
\caption{Stage~3 per-class accuracy on 200 real I/Q captures (50 per class) from the testbed.}
\label{fig:stage3}
\end{figure}

\begin{figure}[t]
\centering
\includegraphics[width=\columnwidth]{fig_spectrograms_best.pdf}
\caption{Spectrograms (224$\times$224) from real I/Q captures on the testbed for four protocol-aware attack types. Each exhibits distinct frequency-time signatures enabling Stage~3 classification.}
\label{fig:spectrograms}
\end{figure}

\subsection{Latency Analysis}

Table~\ref{tab:latency} shows all classification paths remain well within the Near-RT RIC's 100\,ms budget, with the fastest path (ATTACK\_CONFIRMED) completing in under 1\,ms and the heaviest path (FN\_CAUGHT, requiring all three stages) at $\leq$23\,ms (p95). This represents a 7.8$\times$ improvement over Jere et al.'s 180\,ms detection delay~\cite{jere2024}, enabled by the lightweight cascade design that avoids invoking Stage~3 for most samples.

\begin{table}[t]
\centering
\caption{Inference Latency by Classification Path (p5--p95)}
\label{tab:latency}
\small
\begin{tabular}{llr}
\toprule
Path & Stages & Latency (ms) \\
\midrule
CLEAN & S1$\rightarrow$S2 & 0.2--3.0 \\
ATTACK\_CONFIRMED & S1 & $<$0.5 \\
FN\_CAUGHT & S1$\rightarrow$S2$\rightarrow$S3 & 8.8--22.5 \\
PROTOCOL\_AWARE & S1$\rightarrow$S3 & 8.2--20.7 \\
\midrule
\multicolumn{2}{l}{Near-RT RIC Budget} & $\leq$ \textbf{100} \\
\bottomrule
\end{tabular}
\end{table}

\subsection{Comparison with Prior Work}

Table~\ref{tab:comparison} positions MSJC against the most closely related works. Hachimi et al.~\cite{hachimi2020} established the MLP$+$KSVM cascade with four classes (constant, deceptive, random, reactive) but in C-RAN with simulated data and without protocol-specific classification (PSS, PDCCH, DMRS). Jere et al.~\cite{jere2024} validated ML jamming detection on real 5G NR but without O-RAN integration. Moore et al.~\cite{moore2025} and Paz-P\'{e}rez et al.~\cite{pazperez2025} demonstrated O-RAN xApp-based jamming mitigation but with binary detection only. Rahman et al.~\cite{rahman2025} proposed the closed-loop MLOps architecture but limited to binary detection in simulation. MSJC is the first to combine real testbed validation, O-RAN xApp deployment, multi-class classification with protocol-aware sub-typing, sliding-window temporal features, and closed-loop retraining.

\begin{table}[t]
\centering
\caption{Comparison with Related Work}
\label{tab:comparison}
\small
\setlength{\tabcolsep}{2pt}
\begin{tabular}{lcccccc}
\toprule
 & \cite{hachimi2020} & \cite{jere2024} & \cite{moore2025} & \cite{pazperez2025} & \cite{rahman2025} & \textbf{Ours} \\
\midrule
Real testbed & $\times$ & \checkmark & \checkmark & \checkmark & $\times$ & \textbf{\checkmark} \\
O-RAN xApp & $\times$ & $\times$ & \checkmark & \checkmark & \checkmark & \textbf{\checkmark} \\
Multi-class & \checkmark & \checkmark & $\times$ & $\times$ & $\times$ & \textbf{\checkmark} \\
Protocol-aware & $\times$ & $\times$ & $\times$ & $\times$ & $\times$ & \textbf{\checkmark} \\
Sliding window & $\times$ & $\times$ & $\times$ & $\times$ & $\times$ & \textbf{\checkmark} \\
Closed-loop & $\times$ & $\times$ & $\times$ & $\times$ & \checkmark & \textbf{\checkmark} \\
\midrule
Detection & 94.5\% & 98.8\% & 99\% & N/A & N/A & \textbf{98.8\%} \\
False alarm & N/A & N/A & N/A & N/A & N/A & \textbf{5.3\%} \\
\bottomrule
\end{tabular}
\end{table}

\subsection{Closed-Loop MLOps Validation}

We validated the complete closed-loop MLOps cycle on the testbed, extending Rahman et al.'s SAJD architecture~\cite{rahman2025} from binary detection to multi-class classification. The \textbf{Labeler rApp} ingests raw KPI telemetry from InfluxDB in 5-minute sliding windows, applies moving average smoothing, and fits a 2-component GMM for unsupervised Normal/Attack clustering with ARC-based drift detection ($\tau = 0.0004$). The \textbf{Training Manager rApp} monitors xApp accuracy against GMM-generated labels and triggers automated Stage~1 MLP retraining when accuracy drops below 30\%~\cite{rahman2025}. Upon retraining, the Manager deploys the updated model via an A1-like HTTP interface; the xApp performs atomic hot-reload (threading lock $+$ \texttt{state\_dict} swap) without service interruption. In our end-to-end validation, the full cycle---from KPI ingestion through GMM labeling, retraining, and hot-reload---completed within 5~minutes. An InfluxDB~v2 $+$ Grafana observability layer provides real-time monitoring of KPI trends, detection verdicts, and model accuracy.


% ================================================================
\section{Conclusion}
\label{sec:conclusion}

We presented MSJC, a multi-stage jamming classification framework for 5G O-RAN that detects all seven jamming modes---including protocol-aware attacks targeting PSS, PDCCH, and DMRS---with 100\% detection for six of seven attack types (91.9\% for PDCCH), 5.3\% false alarm rate, and $\leq$23\,ms end-to-end latency (p95) on the CCI~xG Testbed, achieving 85\% Stage~3 protocol-aware sub-classification accuracy on real I/Q captures. All results are validated via blocked temporal 5-fold cross-validation with zero sample overlap between train and test sets, a stricter evaluation than the in-sample or hold-out methods used in prior work.

We show that protocol-aware attacks, while invisible to per-sample KPI classifiers, produce distinct multi-dimensional temporal signatures---including subtly depressed SINR and reduced CQI variability---that can be exploited through sliding-window statistical analysis with an RBF-kernel SVM. The three-stage cascade enables cost-effective deployment: most samples are classified in under 1\,ms by Stage~1, with the heavier Stage~3 spectrogram analysis invoked only when protocol-aware attacks are suspected.

Future work includes: (i)~multi-cell extension to detect spatially-distributed jamming, (ii)~active mitigation via E2SM-RC control actions (beam steering, power adaptation)~\cite{abdalla2022}, (iii)~sub-second KPM reporting through gNB scheduler modifications to improve temporal resolution, and (iv)~adversarial robustness against attackers who may adapt to our detection features~\cite{kakani2025}.

\section*{Acknowledgment}
This work was supported by the Commonwealth Cyber Initiative (CCI) and conducted on the CCI xG Testbed~\cite{dasilva2025} at Virginia Tech.

\balance
\bibliographystyle{IEEEtran}
\bibliography{references}

\end{document}
